\subsection{局域辐射源的频谱分析}

前文已从推迟势出发，得到局域电荷体系在长波近似下的远区多极场，以及能量和角动量的瞬时辐射率。这些时域公式给出的是某一时刻源的辐射强度；频谱分析并不引入新的动力学假设，而是把同一组多极矩按时间尺度分解，从而回答辐射能量分布在哪些频率、是否形成连续谱或离散线谱，以及不同频段由轨道的哪些部分控制。因此，具体问题中的势能只负责确定轨道和多极矩，频域中的辐射学部分可以先在一般形式下完成。

以下对非周期过程采用 Fourier 变换约定
\begin{equation*}
	\widetilde f(\omega)=\int_{-\infty}^{\infty}f(t)e^{i\omega t}\d t,
	\qquad
	f(t)=\frac{1}{2\pi}\int_{-\infty}^{\infty}\widetilde f(\omega)e^{-i\omega t}\d\omega.
\end{equation*}
对实函数 $f(t)$ 取共轭即得 $\widetilde f(-\omega)=\widetilde f(\omega)^*$。把 $A(t)$、$B(t)$ 的反变换代入积分并先对时间积分，利用 $\int e^{-i(\omega+\omega')t}\d t=2\pi\delta(\omega+\omega')$，再把积分区间拆成正负频率并两次使用实函数的共轭性质，便得到正频率形式的 Parseval 恒等式
\begin{equation*}
	&\int_{-\infty}^{\infty} A(t):B(t)\d t
	=\frac{1}{2\pi}\int_{-\infty}^{\infty}\widetilde A(\omega):\widetilde B(-\omega)\d\omega\\
	&\quad=\frac{1}{\pi}\int_0^\infty
	\operatorname{Re}\!\left[\widetilde A(\omega):\widetilde B(\omega)^*\right]\d\omega,
\end{equation*}
其中对矢量时冒号理解为普通内积。特别地，分部积分 $n$ 次（非周期过程的相应边界项消失，或可由绝热收敛因子定义）得频率微分规则
\begin{equation*}
	\widetilde{f^{(n)}}(\omega)=(-i\omega)^n\widetilde f(\omega).
\end{equation*}

记观测方向为 $\vb n$，只保留远场中按 $R^{-1}$ 衰减的部分。沿用前文的电偶极矩 $\vb p$、磁偶极矩 $\vb m$ 和无迹电四极矩 $\vb D$，远区电场可统一写成
\begin{equation*}
	\vb E_{\mathrm{rad}}(R\vb n,t)
	=\frac{1}{4\pi\varepsilon_0c^2R}
	\left\{
	\vb n\times\left(\vb n\times\ddot{\vb p}\right)
	+\frac{1}{c}\vb n\times\ddot{\vb m}
	+\frac{1}{6c}\vb n\times
	\left[\vb n\times\left(\dddot{\vb D}\cdot\vb n\right)\right]
	\right\}_{t-R/c}.
\end{equation*}
相应的辐射磁场为 $\vb B_{\mathrm{rad}}=\vb n\times\vb E_{\mathrm{rad}}/c$。对远区电场做 Fourier 变换，迟滞时间换元 $t=t'+R/c$ 只产生相位因子 $e^{i\omega R/c}$，频域远区场为
\begin{equation*}
	\widetilde{\vb E}_{\mathrm{rad}}(R\vb n,\omega)
	=\frac{e^{i\omega R/c}}{4\pi\varepsilon_0c^2R}
	\boldsymbol{\mathcal A}(\omega,\vb n),
\end{equation*}
式中频域辐射振幅定义为
\begin{equation*}
	&\boldsymbol{\mathcal A}(\omega,\vb n)
	=\vb n\times\left(\vb n\times\widetilde{\ddot{\vb p}}\right)
	+\frac{1}{c}\vb n\times\widetilde{\ddot{\vb m}}\\
	&\quad +\frac{1}{6c}\vb n\times
	\left[\vb n\times\left(\widetilde{\dddot{\vb D}}\cdot\vb n\right)\right].
\end{equation*}
辐射磁场使能量以速率 $|\vb E|^2/\mu_0c$ 沿 $\uv n$ 流出，故单位立体角、单位频率的辐射能量为
\begin{equation*}
	\frac{\d^2E}{\d\omega\d\Omega}
	=\frac{R^2}{\pi\mu_0c}
	\left|\widetilde{\vb E}_{\mathrm{rad}}(R\vb n,\omega)\right|^2
	=\frac{1}{16\pi^3\varepsilon_0c^3}
	\left|\boldsymbol{\mathcal A}(\omega,\vb n)\right|^2,
\end{equation*}
其中第二等号是通量对时间的积分与正频率 Parseval 恒等式的直接结果：迟滞相位在模方中抵消，系数合并时用了 $\mu_0\varepsilon_0c^2=1$。上式中的平方必须在各多极振幅相加之后进行，因此在固定观测方向上一般包含 $E1$--$M1$、$E1$--$E2$ 等干涉项。单独取各阶振幅时，有
\begin{equation*}
	\frac{\d^2E_{E1}}{\d\omega\d\Omega}
	=\frac{1}{16\pi^3\varepsilon_0c^3}
	\left|\widetilde{\ddot{\vb p}}\times\vb n\right|^2,
\end{equation*}
\begin{equation*}
	\frac{\d^2E_{M1}}{\d\omega\d\Omega}
	=\frac{1}{16\pi^3\varepsilon_0c^5}
	\left|\widetilde{\ddot{\vb m}}\times\vb n\right|^2,
\end{equation*}
以及
\begin{equation*}
	\frac{\d^2E_{E2}}{\d\omega\d\Omega}
	=\frac{1}{576\pi^3\varepsilon_0c^5}
	\left|\left(\widetilde{\dddot{\vb D}}\cdot\vb n\right)\times\vb n\right|^2.
\end{equation*}
对观测方向积分时使用
\begin{equation*}
	\int n_i n_j\d\Omega=\frac{4\pi}{3}\delta_{ij},
	\qquad
	\int n_i n_jn_kn_l\d\Omega
	=\frac{4\pi}{15}
	\left(\delta_{ij}\delta_{kl}+\delta_{ik}\delta_{jl}+\delta_{il}\delta_{jk}\right).
\end{equation*}
这两个恒等式由各向同性可直接验证：$\int n_in_j\d\Omega$ 只能正比于 $\delta_{ij}$，与 $\delta_{ij}$ 收缩得 $\int\d\Omega=4\pi=3A$；四指标情形按对称性取 $\delta$ 组合的相同系数，收缩两次即定出 $\frac{4\pi}{15}$。以 $E1$ 为例做收缩，$|\widetilde{\ddot{\vb p}}\times\vb n|^2=\widetilde{\ddot p}_i\widetilde{\ddot p}_j^*(\delta_{ij}-n_in_j)$，得 $\int|\widetilde{\ddot{\vb p}}\times\vb n|^2\d\Omega=\frac{8\pi}{3}\widetilde{\ddot p}_i\widetilde{\ddot p}_i^*$；$M1$ 相同。$E2$ 的收缩为 $|(\vb D''\cdot\vb n)\times\vb n|^2=D''_{ij}D''^*_{kl}n_jn_l(\delta_{ik}-n_in_k)$，式中 $\vb D''=\widetilde{\dddot{\vb D}}$，第二个恒等式给出 $\frac{4\pi}{3}-\frac{8\pi}{15}=\frac{4\pi}{5}$ 倍 $D''_{ij}D''^*_{ij}$——无迹条件使 $\delta_{ij}\delta_{kl}$ 项消失。代入各角分布谱并完成系数，任意局域非相对论源的角积分频谱为
\begin{equation*}
	\frac{\d E_{E1}}{\d\omega}
	=\frac{1}{6\pi^2\varepsilon_0c^3}
	\widetilde{\ddot p}_i\widetilde{\ddot p}_i^*
	=\frac{\omega^4}{6\pi^2\varepsilon_0c^3}
	\widetilde p_i\widetilde p_i^*,
\end{equation*}
\begin{equation*}
	\frac{\d E_{M1}}{\d\omega}
	=\frac{1}{6\pi^2\varepsilon_0c^5}
	\widetilde{\ddot m}_i\widetilde{\ddot m}_i^*
	=\frac{\omega^4}{6\pi^2\varepsilon_0c^5}
	\widetilde m_i\widetilde m_i^*,
\end{equation*}
以及
\begin{equation*}
	\frac{\d E_{E2}}{\d\omega}
	=\frac{1}{720\pi^2\varepsilon_0c^5}
	\widetilde{\dddot D}_{ij}\widetilde{\dddot D}_{ij}^*
	=\frac{\omega^6}{720\pi^2\varepsilon_0c^5}
	\widetilde D_{ij}\widetilde D_{ij}^*.
\end{equation*}
第二个等号是频率微分规则的直接结果。不难看到，$E1$、$M1$ 谱按 $\omega^4$ 增长而 $E2$ 谱按 $\omega^6$ 增长，四极项对高频的权重更高。频谱对正频率积分后恢复前文的时域总辐射能量，亦即
\begin{equation*}
	E_X=\int_0^\infty\frac{\d E_X}{\d\omega}\d\omega
	=\int_{-\infty}^{\infty}P_X(t)\d t,
	\qquad X=E1,M1,E2.
\end{equation*}
这个等式既是频谱归一化的依据，也是具体闭式计算最重要的校验之一。

非束缚运动的多极矩在 $t\to\pm\infty$ 具有不同的渐近状态，故上述角积分频谱给出连续谱。在非相对论长波理论内部，最低频项可在求完整 Fourier 积分之前由分部积分的边界项直接确定：例如 $\widetilde{\ddot{\vb p}}(0)=\int\ddot{\vb p}\d t=[\dot{\vb p}]_{-\infty}^{+\infty}$，即
\begin{equation*}
	\widetilde{\ddot{\vb p}}(0)=\dot{\vb p}(+\infty)-\dot{\vb p}(-\infty),
	\qquad
	\widetilde{\ddot{\vb m}}(0)=\dot{\vb m}(+\infty)-\dot{\vb m}(-\infty),
\end{equation*}
\begin{equation*}
	\widetilde{\dddot{\vb D}}(0)=\ddot{\vb D}(+\infty)-\ddot{\vb D}(-\infty).
\end{equation*}
代入前面的角积分频谱公式取零频极限，得到
\begin{equation*}
	&\left.\frac{\d E_{E1}}{\d\omega}\right|_{\omega\to0}
	=\frac{|\Delta\dot{\vb p}|^2}{6\pi^2\varepsilon_0c^3},\\
	&\left.\frac{\d E_{M1}}{\d\omega}\right|_{\omega\to0}
	=\frac{|\Delta\dot{\vb m}|^2}{6\pi^2\varepsilon_0c^5},\\
	&\left.\frac{\d E_{E2}}{\d\omega}\right|_{\omega\to0}
	=\frac{\Delta\ddot D_{ij}\Delta\ddot D_{ij}}{720\pi^2\varepsilon_0c^5}.
\end{equation*}
式中 $\Delta\dot{\vb p}=\dot{\vb p}(+\infty)-\dot{\vb p}(-\infty)$ 为初末渐近值之差，$\Delta\dot{\vb m}$、$\Delta\ddot{\vb D}$ 含义类似。这些式子表明，非相对论多极谱的零频极限只依赖初末渐近运动。然而软辐射的这一性质并不依赖非相对论或长波展开；它实际上是完整 Lienard--Wiechert 理论中的相对论普适结论。下面从点电荷的精确频谱直接推导这一结果，并由此说明其适用范围和更深的红外结构。

考虑若干条点电荷世界线 $\vb r_a(t)$，记
\begin{equation*}
	\boldsymbol\beta_a=\frac{\dot{\vb r}_a}{c},
	\qquad
	\Phi_a(t,\vb n)=t-\frac{\vb n\cdot\vb r_a(t)}{c}.
\end{equation*}
由前文 Lienard--Wiechert 远区场 $\vb E_{\mathrm{rad},a}(R\vb n,t)=\frac{q_a}{4\pi\varepsilon_0cR}\left[\frac{\vb n\times[(\vb n-\boldsymbol\beta_a)\times\dot{\boldsymbol\beta}_a]}{(1-\vb n\cdot\boldsymbol\beta_a)^3}\right]_{t_s}$ 做 Fourier 变换：把观测时间换为粒子固有坐标时，$\d t=(1-\vb n\cdot\boldsymbol\beta_a)\d t_s$，相位为 $\omega\Phi_a$，迟滞因子 $e^{i\omega R/c}$ 在模方中抵消；再按上面的通量--Parseval 程序组装，完整相对论谱角分布为
\begin{equation*}
	\frac{\d^2E}{\d\omega\d\Omega}
	=\frac{1}{16\pi^3\varepsilon_0c}
	\left|
	\sum_a q_a\int_{-\infty}^{\infty}\d t
	\frac{\vb n\times\left[(\vb n-\boldsymbol\beta_a)\times\dot{\boldsymbol\beta}_a\right]}
	{(1-\vb n\cdot\boldsymbol\beta_a)^2}
	e^{i\omega\Phi_a(t,\vb n)}
	\right|^2.
\end{equation*}
这里的时间是粒子世界线的坐标时，式中已经保留了完整的迟滞相位 $\omega\vb n\cdot\vb r/c$ 和全部相对论束射因子。关键恒等式为
\begin{equation*}
	\frac{\d}{\d t}
	\left[
	\frac{\vb n\times(\vb n\times\boldsymbol\beta_a)}
	{1-\vb n\cdot\boldsymbol\beta_a}
	\right]
	=
	\frac{\vb n\times\left[(\vb n-\boldsymbol\beta_a)\times\dot{\boldsymbol\beta}_a\right]}
	{(1-\vb n\cdot\boldsymbol\beta_a)^2}.
\end{equation*}
它可由三重矢量积展开 $\vb n\times(\vb n\times\boldsymbol\beta)=\vb n(\vb n\cdot\boldsymbol\beta)-\boldsymbol\beta$ 与商的求导直接验证：分子展开为 $(\vb n\cdot\dot{\boldsymbol\beta})\vb n-(\vb n\cdot\dot{\boldsymbol\beta})\boldsymbol\beta-(1-\vb n\cdot\boldsymbol\beta)\dot{\boldsymbol\beta}$，恰为 $\vb n\times[(\vb n-\boldsymbol\beta)\times\dot{\boldsymbol\beta}]$ 的三重矢量积展开。若各粒子在 $t\to\pm\infty$ 趋于匀速，便可在精确谱角分布中先取 $\omega\to0$：此时相位 $e^{i\omega\Phi_a}\to1$，积分即全导数的边界值。将所有渐近外线统一编号为 $A$，约定出射线 $\eta_A=+1$、入射线 $\eta_A=-1$，得到领先软振幅
\begin{equation*}
	\boldsymbol{\mathcal S}^{(0)}(\vb n)
	=\sum_A\eta_Aq_A
	\frac{\vb n\times(\vb n\times\boldsymbol\beta_A)}
	{1-\vb n\cdot\boldsymbol\beta_A}.
\end{equation*}
于是任意相对论散射的领先软谱为
\begin{equation*}
	\left.\frac{\d^2E}{\d\omega\d\Omega}\right|_{\omega\to0}
	=\frac{1}{16\pi^3\varepsilon_0c}
	\left|\boldsymbol{\mathcal S}^{(0)}(\vb n)\right|^2.
\end{equation*}
这个结果只含各带电粒子的初末速度，不含相互作用区内的加速度历史。换言之，求领先软辐射时不需要先求完整的时域辐射波形；动力学只需提供渐近散射映射。

领先软谱还可以完成全立体角积分。物理散射过程满足电荷守恒
\begin{equation*}
	\sum_A\eta_Aq_A=0.
\end{equation*}
展开模方得 $\sum_{A,B}\eta_A\eta_Bq_Aq_B\frac{g(\boldsymbol\beta_A)\cdot g(\boldsymbol\beta_B)}{(1-\vb n\cdot\boldsymbol\beta_A)(1-\vb n\cdot\boldsymbol\beta_B)}$，其中 $g(\boldsymbol\beta)=\vb n\times(\vb n\times\boldsymbol\beta)$，且
\begin{equation*}
	g(\boldsymbol\beta_A)\cdot g(\boldsymbol\beta_B)
	=\boldsymbol\beta_A\cdot\boldsymbol\beta_B-(\vb n\cdot\boldsymbol\beta_A)(\vb n\cdot\boldsymbol\beta_B).
\end{equation*}
把 $(\vb n\cdot\boldsymbol\beta_A)(\vb n\cdot\boldsymbol\beta_B)$ 拆成 $1-(1-\vb n\cdot\boldsymbol\beta_A)-(1-\vb n\cdot\boldsymbol\beta_B)+(1-\vb n\cdot\boldsymbol\beta_A)(1-\vb n\cdot\boldsymbol\beta_B)$，除以分母后，常数项 $-(\sum_A\eta_Aq_A)^2$ 与含 $(1-\vb n\cdot\boldsymbol\beta_A)^{-1}$ 的单项（因子 $\sum_B\eta_Bq_B=0$）均因电荷守恒为零，只余
\begin{equation*}
	\left|\boldsymbol{\mathcal S}^{(0)}(\vb n)\right|^2
	=-
	\sum_{A,B}\eta_A\eta_Bq_Aq_B
	\frac{1-\boldsymbol\beta_A\cdot\boldsymbol\beta_B}
	{(1-\vb n\cdot\boldsymbol\beta_A)(1-\vb n\cdot\boldsymbol\beta_B)}.
\end{equation*}
所需角积分可以从两个初等公式推出：
\begin{equation*}
	\frac{1}{AB}=\int_0^1\frac{\d x}{[xA+(1-x)B]^2},
	\qquad
	\int\frac{\d\Omega}{(1-\vb n\cdot\vb w)^2}
	=\frac{4\pi}{1-w^2},
	\quad w<1.
\end{equation*}
第二个等式取极轴沿 $\vb w$ 后 $2\pi\int_{-1}^{1}\frac{\d c}{(1-wc)^2}$ 可直接积出。令
\begin{equation*}
	\Gamma_{AB}
	=\gamma_A\gamma_B
	\left(1-\boldsymbol\beta_A\cdot\boldsymbol\beta_B\right),
	\qquad
	\beta_{AB}=\sqrt{1-\Gamma_{AB}^{-2}},
\end{equation*}
其中 $\Gamma_{AB}$ 是两条渐近外线之间的相对 Lorentz 因子，$\beta_{AB}$ 是相应的相对速度。对成对求和形式使用 Feynman 参数，把两个分母合并在 $\vb w(x)=x\boldsymbol\beta_A+(1-x)\boldsymbol\beta_B$ 中并交换积分次序，得
\begin{equation*}
	\int\d\Omega\,
	\frac{1-\boldsymbol\beta_A\cdot\boldsymbol\beta_B}
	{(1-\vb n\cdot\boldsymbol\beta_A)(1-\vb n\cdot\boldsymbol\beta_B)}
	=4\pi
	(1-\boldsymbol\beta_A\cdot\boldsymbol\beta_B)
	\int_0^1\frac{\d x}
	{1-|x\boldsymbol\beta_A+(1-x)\boldsymbol\beta_B|^2}.
\end{equation*}
余下的 $x$ 积分是配平方后的初等积分：分母 $1-|\vb w(x)|^2$ 为 $x$ 的二次式，平移区间后化为 $\int_{-1}^{1}\frac{\d u}{\alpha^2-\gamma^2u^2}$ 型，由 $\Gamma_{AB}$ 的定义化简为
\begin{equation*}
	\int_0^1\frac{\d x}{1-|x\boldsymbol\beta_A+(1-x)\boldsymbol\beta_B|^2}
	=\frac{\operatorname{artanh}\beta_{AB}}
	{\beta_{AB}(1-\boldsymbol\beta_A\cdot\boldsymbol\beta_B)}.
\end{equation*}
（等速率、夹角 $\vartheta$ 时可直接复核：$1-|\vb w|^2=1-\beta^2\cos^2\frac{\vartheta}{2}-\beta^2\sin^2\frac{\vartheta}{2}(2x-1)^2$，平移区间后即化为 $\frac{1}{ab}\operatorname{artanh}(b/a)$ 型初等积分。）因此角积分后的相对论软谱具有完全闭合的形式
\begin{equation*}
	\left.\frac{\d E}{\d\omega}\right|_{\omega\to0}
	=-\frac{1}{4\pi^2\varepsilon_0c}
	\sum_{A,B}\eta_A\eta_Bq_Aq_B
	\left[
	\frac{\operatorname{artanh}\beta_{AB}}{\beta_{AB}}-1
	\right].
\end{equation*}
方括号中的 $-1$ 可以加入，是因为电荷守恒条件使相应的双重求和为零。当 $A=B$ 时，按连续极限 $\operatorname{artanh}\beta_{AA}/\beta_{AA}\to1$，自项自动消失。

最简单而重要的情形是一枚电荷 $q$ 由初速度 $\boldsymbol\beta_i$ 偏转到末速度 $\boldsymbol\beta_f$。此时只有一对入射、出射外线，闭合软谱中非零的只有 $A\neq B$ 两项，化为
\begin{equation*}
	\left.\frac{\d E}{\d\omega}\right|_{\omega\to0}
	=\frac{q^2}{2\pi^2\varepsilon_0c}
	\left[
	\frac{\operatorname{artanh}\beta_{if}}{\beta_{if}}-1
	\right]
	=\frac{q^2}{2\pi^2\varepsilon_0c}
	\left[
	\frac{\Gamma_{if}\operatorname{arccosh}\Gamma_{if}}
	{\sqrt{\Gamma_{if}^2-1}}-1
	\right].
\end{equation*}
若初末速率相同为 $\beta c$，夹角为 $\Theta$，则
\begin{equation*}
	\Gamma_{if}
	=\gamma^2(1-\beta^2\cos\Theta)
	=1+2\gamma^2\beta^2\sin^2\frac{\Theta}{2},
\end{equation*}
第二个等号用了 $\gamma^{-2}=1-\beta^2$ 与倍角公式。在非相对论极限，$\beta_{if}=|\boldsymbol\beta_f-\boldsymbol\beta_i|+O(\beta^3)$——把 $\beta_{if}^2=1-\Gamma_{if}^{-2}$ 按 $\beta^2$ 展开得 $\frac{(1-\beta^2\cos\Theta)^2-(1-\beta^2)^2}{(1-\beta^2\cos\Theta)^2}\simeq4\beta^2\sin^2\frac{\Theta}{2}$——并且
\begin{equation*}
	\frac{\operatorname{artanh}x}{x}-1
	=\frac{x^2}{3}+\frac{x^4}{5}+O(x^6),
\end{equation*}
取自 $\operatorname{artanh}$ 的幂级数。故单电荷软谱恢复
\begin{equation*}
	\left.\frac{\d E}{\d\omega}\right|_{\omega\to0}
	=\frac{q^2|\Delta\vb v|^2}{6\pi^2\varepsilon_0c^3}
	+O(c^{-5}),
\end{equation*}
这正是前面非相对论边界形式中的电偶极部分。相反，当 $\Gamma_{if}\gg1$ 时，
\begin{equation*}
	\frac{\Gamma_{if}\operatorname{arccosh}\Gamma_{if}}
	{\sqrt{\Gamma_{if}^2-1}}-1
	=\ln(2\Gamma_{if})-1+O(\Gamma_{if}^{-2}\ln\Gamma_{if}),
\end{equation*}
由 $\operatorname{arccosh}\Gamma\simeq\ln(2\Gamma)$ 展开得到。软能谱因相对论前向束射而出现对数增强。若散射角很小，则
\begin{equation*}
	\Gamma_{if}-1\simeq\frac{1}{2}\gamma^2\Theta^2.
\end{equation*}
于是 $\gamma\Theta\ll1$ 时软谱正比于 $\gamma^2\Theta^2$，而 $\gamma\Theta\gg1$ 时则按 $2\ln(\gamma\Theta)$ 增长。

零频振幅还等于远区辐射场对迟滞时间的积分。频域远区场在零频处的值正是该积分：$\widetilde{\vb E}_{\mathrm{rad}}(R\vb n,0)=\int\vb E_{\mathrm{rad}}(R\vb n,u)\d u$；沿用上面的边界值推导，若 $u=t_{\mathrm{obs}}-R/c$，则
\begin{equation*}
	\int_{-\infty}^{\infty}\vb E_{\mathrm{rad}}(R\vb n,u)\d u
	=\frac{1}{4\pi\varepsilon_0cR}
	\boldsymbol{\mathcal S}^{(0)}(\vb n).
\end{equation*}
因此领先软定理与电磁记忆效应是同一个边界量的频域和时域表述。另一方面，由于
\begin{equation*}
	\frac{\d N_\gamma}{\d\omega}
	=\frac{1}{\hbar\omega}\frac{\d E}{\d\omega},
\end{equation*}
频率为 $\omega$ 的光子携带能量 $\hbar\omega$。当 $\d E/\d\omega$ 在 $\omega\to0$ 趋于非零常数时，软光子数按 $1/\omega$ 发散，而总软辐射能量 $\int_0^{\omega_{\max}}(\d E/\d\omega)\d\omega$ 仍然有限。实验中有限观测时间、能量分辨率以及虚光子修正共同决定可分辨的红外下限。

还需注意，长程相互作用的次领先软展开一般不是普通的 $\omega$ 幂级数。短程力下，轨道在足够大的 $|t|$ 上严格趋于直线，领先软项之后可按 $\omega$ 展开；库仑力则给出
\begin{equation*}
	\vb r(t)=c\boldsymbol\beta_\pm t+\vb c_\pm\ln|t|+\vb b_\pm+O(t^{-1}),
	\qquad
	\boldsymbol\beta(t)=\boldsymbol\beta_\pm+\frac{\vb c_\pm}{ct}+O(t^{-2}).
\end{equation*}
其本质可从尾积分看出。对 $a>0$，由分部积分或指数积分定义，
\begin{equation*}
	\int_T^\infty\frac{e^{ia\omega t}}{t^2}\d t
	=\frac{e^{ia\omega T}}{T}+ia\omega E_1(-ia\omega T)
	=\frac{1}{T}-ia\omega\ln(\omega T)+O(\omega),
\end{equation*}
其中 $E_1$ 为指数积分，最后一个等号只显示了非解析部分。故四维库仑散射的软振幅一般具有
\begin{equation*}
	\boldsymbol{\mathcal A}(\omega,\vb n)
	=\boldsymbol{\mathcal S}^{(0)}(\vb n)
	+\omega\ln\omega\,\boldsymbol{\mathcal S}^{(\ln)}(\vb n)
	+O(\omega),
\end{equation*}
而不是简单的 $\boldsymbol{\mathcal S}^{(0)}+\omega\boldsymbol{\mathcal S}^{(1)}+\cdots$。领先常数项仍由领先软因子完全确定，但若要求次领先精度，就必须保留库仑轨道的长程对数尾。

硬频渐近则由时间复平面中离实轴最近的奇点或转折点控制。若某一实际进入辐射振幅的函数在上半平面最近奇点 $t_\star=t_0+i\tau_\star$ 附近具有
\begin{equation*}
	F(t)\sim A(t-t_\star)^{-\lambda},
	\qquad \tau_\star>0,
\end{equation*}
则由鞍点法（陡降路径）的初等形式，指数由奇点到实轴的距离决定、幂由奇点的指数决定，
\begin{equation*}
	\widetilde F(\omega)
	\sim C\,\omega^{\lambda-1}e^{i\omega t_\star},
	\qquad
	|\widetilde F(\omega)|^2
	\sim |C|^2\omega^{2\lambda-2}e^{-2\omega\tau_\star}.
\end{equation*}
因此硬频指数截止给出最近复时间奇点到实轴的距离；当奇点随轨道参数趋近实轴，或积分端点本身不光滑时，指数尾会退化为幂律或转折点一致渐近。

若运动是周期为 $T=2\pi/\Omega$ 的束缚运动，应使用 Fourier 级数
\begin{equation*}
	f_N=\frac{1}{T}\int_0^T f(t)e^{iN\Omega t}\d t,
	\qquad
	f(t)=\sum_{N=-\infty}^{\infty}f_Ne^{-iN\Omega t}.
\end{equation*}
实函数满足 $f_{-N}=f_N^*$。把连续谱公式中的频域多极矩换成周期情形的 $\widetilde{\vb p}(\omega)=2\pi\sum_N\vb p_N\delta(\omega-N\Omega)$，对正频率积分时只有 $N\geq1$ 的谱线贡献，例如
\begin{equation*}
	E_{E1}=\frac{1}{6\pi^2\varepsilon_0c^3}
	\int_0^\infty\omega^4|\widetilde{\vb p}(\omega)|^2\d\omega
	=\frac{4\pi^2}{6\pi^2\varepsilon_0c^3}
	\sum_{N=1}^{\infty}(N\Omega)^4p_{i,N}p_{i,N}^*,
\end{equation*}
其中利用了 $\delta$ 函数在正半轴的分离性 $\int_0^\infty\delta(\omega-N\Omega)\delta(\omega-M\Omega)\d\omega=\delta_{NM}$。每个周期的平均功率为 $\overline P=E/T=E\Omega/2\pi$，故第 $N\geq1$ 条谱线的周期平均功率为
\begin{equation*}
	&P_{E1,N}=\frac{(N\Omega)^4}{3\pi\varepsilon_0c^3}p_{i,N}p_{i,N}^*,\\
	&P_{M1,N}=\frac{(N\Omega)^4}{3\pi\varepsilon_0c^5}m_{i,N}m_{i,N}^*,\\
	&P_{E2,N}=\frac{(N\Omega)^6}{360\pi\varepsilon_0c^5}D_{ij,N}D_{ij,N}^*.
\end{equation*}
$M1$、$E2$ 由同样的代换得到。所有正整数谐波求和后恢复周期平均的时域功率：
\begin{equation*}
	\overline P_X=\sum_{N=1}^{\infty}P_{X,N}.
\end{equation*}
这可由周期 Parseval 恒等式 $\int_0^T|f(t)|^2\d t=T\sum_N|f_N|^2$ 直接验证。由此可见，连续谱和离散线谱只是 Fourier 变换与 Fourier 级数在非周期、周期运动中的两种表现：周期运动的频域场是 $\delta$ 函数谱，$\delta$ 平方在连续谱归一化中没有定义，必须改用周期 Parseval——这正是两者不能混用归一化的原因。下面以圆周运动为例，说明周期谱线公式在精确相对论运动中的用法。

\begin{example}
圆周运动的同步辐射谱.半径为 $a$、速率为 $\beta c$ 的圆周运动可以完全积分，是检验周期谱线公式并展示相对论效应如何改造谱形的最佳范例。取轨道位于 $xy$ 平面内
\begin{equation*}
	\vb r_q(t)=a(\sin\varphi,-\cos\varphi,0),
	\qquad
	\boldsymbol\beta=\beta(\cos\varphi,\sin\varphi,0),
	\qquad
	\varphi=\Omega t,
	\qquad
	\Omega=\frac{\beta c}{a}.
\end{equation*}
取场点方向与偏振基矢为
\begin{equation*}
	\uv R=(\cos\theta,0,\sin\theta),
	\qquad
	\uv e_a=(-\sin\theta,0,\cos\theta),
	\qquad
	\uv e_b=(0,1,0).
\end{equation*}

谐波谱的由来。由于轨道以 $T=2\pi/\Omega$ 为周期，精确谱积分中的被积函数（含完整迟滞相位 $e^{-i\omega\uv R\cdot\vb r_q/c}$）是固有时间 $t_s$ 的周期函数，其 Fourier 变换化为一列 $\delta$ 谱线：对任意周期函数 $g$，
\begin{equation*}
	\int_{-\infty}^{\infty}g(t_s)e^{i\omega t_s}\d t_s
	=2\pi\sum_{m=-\infty}^{\infty}g_m\delta(\omega-m\Omega),
	\qquad
	g_m=\frac{1}{T}\int_0^Tg(t_s)e^{im\Omega t_s}\d t_s.
\end{equation*}
这就是频率 $m\Omega=m\beta c/a$ 处谱线的来源。先用前面的全导数恒等式对精确谱积分分部积分（周期运动的边界项消失），横向矢量化为 $\uv R\times(\uv R\times\boldsymbol\beta)$；在 $(\uv e_a,\uv e_b,\uv R)$ 基中
\begin{equation*}
	\uv R\times(\uv R\times\boldsymbol\beta)
	=\beta\sin\theta\cos\varphi\,\uv e_a-\beta\sin\varphi\,\uv e_b,
	\qquad
	e^{-i\omega\uv R\cdot\vb r_q/c}=e^{-iu\sin\varphi},
	\quad u=\frac{\omega a\cos\theta}{c}.
\end{equation*}
对 $e^{-iu\sin\varphi}$ 用 Jacobi--Anger 展开 $e^{-iu\sin\varphi}=\sum_m J_m(u)e^{-im\varphi}$：$\cos\varphi\,e^{-iu\sin\varphi}$ 的 $m$ 阶 Fourier 系数为 $\frac12[J_{m-1}(u)+J_{m+1}(u)]=\frac{m}{u}J_m(u)$，$\sin\varphi\,e^{-iu\sin\varphi}$ 的为 $\frac{1}{2i}[J_{m+1}(u)-J_{m-1}(u)]=iJ_m'(u)$，其中用了 Bessel 递推。完成系数后，远区电场为
\begin{equation*}
	&\vb E_{\mathrm{rad}}(\vb r,\omega)
	=-\frac{iq\omega\beta e^{i\omega r/c}}{2\varepsilon_0cr}
	\sum_{m=-\infty}^{\infty}
	\delta\left(\omega-m\frac{\beta c}{a}\right)\\
	&\quad\times
	\left[
	-i\uv e_b J_m'\left(\frac{\omega a\cos\theta}{c}\right)
	+\uv e_a m\sin\theta
	\frac{J_m\left(\omega a\cos\theta/c\right)}{\omega a\cos\theta/c}
	\right].
\end{equation*}
故频谱由频率 $\omega=m\beta c/a$ 处的 $\delta$ 函数谱线构成：$m$ 次谐波中垂直于轨道平面的偏振分量由 $J_m'$ 控制，面内偏振分量由 $m\sin\theta\,J_m/(m\beta\cos\theta)$ 控制。对 $m\geq1$ 谐波应用周期 Parseval 恒等式，每周期能量在单位立体角中的分布为
\begin{equation*}
	&\frac{\d W_m}{\d\Omega}
	=\frac{2r^2T}{\mu_0c}|\vb E_m|^2
	=\frac{Tq^2\omega^2\beta^2}{8\pi^2\varepsilon_0c}
	\left[J_m'^2(m\beta\cos\theta)\right.\\
	&\qquad
	\left.+\frac{\sin^2\theta}{\beta^2\cos^2\theta}J_m^2(m\beta\cos\theta)\right],
\end{equation*}
其中 $\vb E_m=\widetilde{\vb E}_m/2\pi$，正负谐波贡献相同，最后在 $\omega=m\Omega$、$u=m\beta\cos\theta$ 处取值并完成系数。对 $m\geq1$ 的谐波求和可封闭完成，得到每周期能量的角分布
\begin{equation*}
	\frac{\d W}{\d\Omega}
	=\frac{q^2\beta^3}{64\pi\varepsilon_0a}
	\frac{4(1+\sin^2\theta)-\beta^2\cos^4\theta(1+3\beta^2)}
	{(1-\beta^2\cos^2\theta)^{7/2}},
\end{equation*}
再对立体角积分得每周期总能量
\begin{equation*}
	W=2\pi\int_{-\pi/2}^{\pi/2}\frac{\d W}{\d\Omega}\cos\theta\d\theta
	=\frac{q^2\gamma^4\beta^3}{3\varepsilon_0a},
\end{equation*}
折回功率 $P=W/T=q^2\gamma^4\beta^4c/(6\pi\varepsilon_0a^2)$，与相对论 Larmor 公式一致，验证了谱线求和校验在精确相对论情形下仍然成立。

连续化极限。当速度接近光速 ($\gamma\gg1$)、观测角很小 ($\theta\ll1$) 时，谱的特征宽度 $\sim\gamma^3\Omega$ 远大于谐波间距 $\Omega$，每个特征宽度内包含大量谐波，求和可化为积分：对缓慢变化的 $f$，
\begin{equation*}
	\sum_{m=-\infty}^{\infty}f(m)\delta(\omega-m\Omega)
	=\frac{1}{\Omega}f\left(\frac{\omega}{\Omega}\right),
	\qquad
	\Omega=\frac{\beta c}{a}.
\end{equation*}
此时谐波参数 $u=m\beta\cos\theta\simeq m\left(1-\frac{\theta^2+\gamma^{-2}}{2}\right)$ 与阶数 $m$ 之差为 $\frac{m}{2}(\theta^2+\gamma^{-2})$ 量级，适用 Bessel 函数的 Debye 一致渐近（可由 $K_\nu$ 的积分表示 $K_\nu(z)=\int_0^\infty e^{-z\cosh t}\cosh(\nu t)\d t$ 复核），上述线性组合化为修正贝塞尔函数
\begin{equation*}
	&\sum_{m=-\infty}^{\infty}
	\delta\left(\omega-m\frac{\beta c}{a}\right)
	\left(
	-i\uv e_b J_m'\left(\frac{\omega a\cos\theta}{c}\right)
	+\uv e_a\sin\theta
	\frac{mJ_m\left(\omega a\cos\theta/c\right)}{\omega a\cos\theta/c}
	\right)\\
	&\quad=-\frac{ia}{\sqrt{3}\pi\beta c}(\theta^2+\gamma^{-2})
	\left(
	\uv e_b K_{2/3}(\xi)
	+i\uv e_a\frac{\theta}{\sqrt{\theta^2+\gamma^{-2}}}K_{1/3}(\xi)
	\right),
	\qquad
	\xi=\frac{\omega a}{3c}(\theta^2+\gamma^{-2})^{3/2},
\end{equation*}
其中 $\xi=\frac{m}{3}(\theta^2+\gamma^{-2})^{3/2}$ 在 $\omega=m\Omega$ 处的值，且 $\beta\simeq1$。代入每周期能量谱即得连续化的角分布频谱
\begin{equation*}
	&\frac{\d^2W}{\d\Omega\d\omega}
	=\frac{q^2}{12\pi^3\varepsilon_0c}
	\left(\frac{\omega a}{c}\right)^2(\theta^2+\gamma^{-2})^2
	\left(
	(K_{2/3}(\xi))^2
	+\frac{\theta^2}{\theta^2+\gamma^{-2}}(K_{1/3}(\xi))^2
	\right),\\
	&\frac{\d W}{\d\Omega}
	=\frac{7q^2}{64\pi\varepsilon_0a(\theta^2+\gamma^{-2})^{5/2}}
	\left(
	1+\frac{5}{7}\frac{\theta^2}{\theta^2+\gamma^{-2}}
	\right),\\
	&W=2\pi\int_{-\infty}^{\infty}\frac{\d W}{\d\Omega}\d\theta
	=\frac{q^2\gamma^4}{3\varepsilon_0a}.
\end{equation*}
修正贝塞尔函数在 $\xi\gg1$ 时按 $K_\nu(\xi)\sim\sqrt{\pi/2\xi}\,e^{-\xi}$ 指数衰减，故每个观测方向上辐射显著的频率满足 $\xi\lesssim1$，即 $\theta^2+\gamma^{-2}\lesssim(3c/\omega a)^{2/3}$；在谱峰频率 $\omega\sim\gamma^3\beta c/a$ 处——$\theta=0$ 时谱 $\propto z^2K_{2/3}^2(z)$，$z=\omega a\gamma^{-3}/3c$，而 $z^2K_{2/3}^2(z)$ 在 $z\sim1$ 处最大——这一判据给出 $\theta\lesssim1/\gamma$。谱集中在前向 $\theta\sim1/\gamma$ 的窄锥内并展宽到约 $\gamma^3$ 条谐波：低速时圆周运动只辐射 $m=1$ 基波，$\gamma\gg1$ 时谱被相对论束射聚束展宽——相对论效应正是以这种锥形结构改造了谱形。
\end{example}

在经典散射问题中，固定碰撞参数 $\vb b$ 的 $\d E/\d\omega$ 描述一条给定轨道的辐射。若入射束在碰撞参数平面上均匀（单位面积、单位频率的注入能量相同），则定义能量加权的谱辐射截面
\begin{equation*}
	\frac{\d\chi_X}{\d\omega}
	=\int\d^2b\,\frac{\d E_X(\vb b)}{\d\omega}.
\end{equation*}
轴对称散射中 $\d^2b=b\d b\d\phi$ 的角向积分给 $2\pi$，即 $\d^2b=2\pi b\d b$。频率为 $\omega$ 的光子携带能量 $\hbar\omega$，相应的光子数截面为
\begin{equation*}
	\frac{\d\sigma_{\gamma,X}}{\d\omega}
	=\frac{1}{\hbar\omega}\frac{\d\chi_X}{\d\omega}.
\end{equation*}
谱辐射截面的大碰撞参数端通常控制红外行为，小碰撞参数端通常控制硬频行为；屏蔽、有限尺寸、量子反冲和相对论效应分别为这些端点提供物理截断。

前面的角积分频谱及周期谱线公式属于非相对论长波多极理论，其共同条件是
\begin{equation*}
	\frac{v_{\mathrm{src}}}{c}\ll1,
	\qquad
	\frac{\omega a_{\mathrm{eff}}(\omega)}{c}\ll1.
\end{equation*}
后一条件必须对实际贡献于该频率的轨道区段检查，而不能只以无穷远速度或某个固定几何尺度代替。通常 $E2$ 振幅相对 $E1$ 振幅低一阶 $\omega a/c$；若电偶极矩因对称性或荷质比关系抵消，则必须把四极项作为领先项保留。与此不同，领先软谱角分布和角积分软谱是完整相对论电动力学在 $\omega\to0$ 的领先结论，不要求 $v/c\ll1$ 或长波多极展开；它们的限制是渐近入射、出射状态必须存在，并且所研究频率满足软条件 $\omega\tau_{\mathrm{sc}}\ll1$。

因此，一个具体辐射问题的频谱分析应先判断所需的是完整频谱还是仅需软区。完整的非相对论频谱应先由动力学求出轨道并构造 $\vb p(t)$、$\vb m(t)$、$\vb D(t)$，再根据运动是否周期选择连续 Fourier 变换或 Fourier 级数约定；所得频域多极矩代入完整角分布或相应的角积分频谱公式。若只研究任意速度散射的领先软辐射，则不必求完整轨道 Fourier 变换，只需由散射动力学确定渐近四动量，再代入领先软谱角分布或角积分软谱。散射系综最后按谱辐射截面定义对碰撞参数积分；频谱归一化恒等式、非相对论软极限公式、相对论软极限公式和硬频渐近公式分别检查归一化与三个渐近区的行为。下面把这些一般性结论用于具体的运动模型，依次展示完整非相对论频谱与相对论固定中心软谱的计算。
